\label{sec:appendix}
\section{Dataset Details}
\label{appendix:dataset}


As discussed in~\refsec{sec:task_dataset}, we curate the FreshWiki dataset by collecting \textit{recent} and \textit{high-quality} English Wikipedia articles. We select the most-edited pages over a specific period rather than using creation dates as a cutoff because most of Wikipedia articles are ``stubs'' or are of low quality when they were created. For quality, we consider articles predicted to be of B-class quality or above. According to Wikipedia statistics\footnote{\url{https://en.wikipedia.org/wiki/Wikipedia:Content_assessment}}, only around 3\% of existing Wikipedia pages meet this quality standard. As LLMs can generate reasonably good outputs, we think it is important to use high-quality human-written articles as references for further research.

\begin{table}
\centering
\resizebox{0.85\columnwidth}{!}{%
\begin{tabular}{lc} 
\toprule
Average Numer of Sections            & 8.4     \\
Average Number of All-level Headings & 15.8    \\ 
\midrule
Average Length of a Section          & 327.8   \\
Average Length of Total Article      & 2159.1  \\ 
\midrule
Average Number of References         & 90.1    \\
\bottomrule
\end{tabular}
}
\caption{Statistics of the dataset used in our experiments.}
\label{table:data_statistics}
\end{table}

\begin{figure}[t]
    \centering 
    \resizebox{\columnwidth}{!}{%
    \includegraphics{img/ref_trend.png}
    }
    \caption{Evolution of reference count in the Wikipedia article editing process.
    }
    \label{Fig.ref_trend}
\end{figure}

\begin{figure}[t]
    \centering 
    \resizebox{\columnwidth}{!}{%
    \includegraphics{img/edit_distribution.png}
    }
    \caption{Distribution of edit counts for Wikipedia articles in our experiments ($n=100$). 
    }
    \label{Fig.edit_distribution}
\end{figure}

For experiments in this work, we randomly select 100 samples with human-written articles under 3000 words to have a meaningful comparison. \reftab{table:data_statistics} gives the data statistics. Notably, human-authored articles have a large number of references but they require numerous edits to achieve this. \reffig{Fig.ref_trend} illustrates the evolution of the reference count in the article edit process and \reffig{Fig.edit_distribution} gives the distribution of edit counts for human-authored articles used in our experiments.


\section{Pseudo Code of \system}
\label{appendix:pseudo_code}
In~\refsec{sec:method}, we introduce \system, a framework that automates the pre-writing stage by discovering different perspectives, simulating information-seeking conversations, and creating a comprehensive outline. Algorithm~\ref{code:pipeline} displays the skeleton of \system.

We implement \system with zero-shot prompting using the DSPy framework~\citep{khattab2023dspy}. Listing~\ref{lst:prompt1} and~\ref{lst:prompt2} show the prompts used in our implementation. We highlight that \system offers a general framework designed to assist the creation of grounded, long-form articles, without depending extensively on prompt engineering for a single domain.

\begin{algorithm}
    \SetKwFunction{findRelatedTopics}{gen\_related\_topics}
    \SetKwFunction{getWikiArticle}{get\_wiki\_article}
    \SetKwFunction{extractTOC}{extract\_toc}
    \SetKwFunction{identifyPerspectives}{gen\_perspectives}
    \SetKwFunction{askQuestion}{gen\_qn}
    \SetKwFunction{questionToQuery}{gen\_queries}
    \SetKwFunction{search}{search\_and\_sift}
    \SetKwFunction{answer}{gen\_ans}
    \SetKwFunction{genOutline}{direct\_gen\_outline}
    \SetKwFunction{refineOutline}{refine\_outline}
    \SetKwInOut{KwIn}{Input}
    \SetKwInOut{KwOut}{Output}

    \KwIn{Topic $t$, maximum perspective $N$, maximum conversation round $M$}
    \KwOut{Outline $\mathcal{O}$, references $\mathcal{R}$}

    P0 = "basic fact writer ..." \tcp{Constant.}

    $\mathcal{R} \leftarrow [\ ]$

    \tcp{Discover perspectives $\mathcal{P}$.}
    
    related\_topics $\leftarrow$ $\findRelatedTopics$($t$)

    tocs $\leftarrow [\ ]$ 

    \ForEach{\upshape{related\_t} in related\_topics}{
        article $\leftarrow$ $\getWikiArticle$(related\_t)
        \If{\upshape{article}}{
            tocs.append($\extractTOC$(article))
        }
    }

    $\mathcal{P} \leftarrow$ $\identifyPerspectives$($t$, tocs)

    $\mathcal{P} \leftarrow$ [P0] + $\mathcal{P}$[:$N$]

    \tcp{Simulate conversations.}
    convos $\leftarrow$ $[\ ]$
    
    \ForEach{$p$ in $\mathcal{P}$}{
        convo\_history $\leftarrow$ $[\ ]$
        
        \For{$i$ = $1$ \upshape{to} $M$}{
            \tcp{Question asking.}
            
            $q$ $\leftarrow$ $\askQuestion$($t$, $p$, dlg\_history)
            
            convo\_history.append($q$)

            \tcp{Question answering.}

            queries $\leftarrow$ $\questionToQuery$($t$, q)

            sources $\leftarrow$ $\search$(queries)

            $a$ $\leftarrow$ $\answer$($t$, q, sources)

            convo\_history.append($a$)

            $\mathcal{R}$.append(sources)
        }

        convos.append(convo\_history)
        
        
    }

    \tcp{Create the outline.}

    $\mathcal{O}_\text{D}$ $\leftarrow$ $\genOutline$($t$)

    $\mathcal{O}$ $\leftarrow$ $\refineOutline$($t$, $\mathcal{O}_\text{D}$, convos)
    
    \KwRet{$\mathcal{O}$, $\mathcal{R}$}

    \caption{\system}
    \label{code:pipeline}
\end{algorithm}

\begin{figure*}
    \begin{lstlisting}[language=Python]
class GenRelatedTopicsPrompt(dspy.Signature):
    """
    I'm writing a Wikipedia page for a topic mentioned below. Please identify and recommend some Wikipedia pages on closely related subjects. I'm looking for examples that provide insights into interesting aspects commonly associated with this topic, or examples that help me understand the typical content and structure included in Wikipedia pages for similar topics.
    Please list the urls in separate lines.
    """

    topic = dspy.InputField(prefix="Topic of interest:", format=str)
    related_topics = dspy.OutputField()

class GenPerspectivesPrompt(dspy.Signature):
    """
    You need to select a group of Wikipedia editors who will work together to create a comprehensive article on the topic. Each of them represents a different perspective, role, or affiliation related to this topic. You can use other Wikipedia pages of related topics for inspiration. For each editor, add description of what they will focus on.
    Give your answer in the following format: 1. short summary of editor 1: description\n2. short summary of editor 2: description\n...
    """

    topic = dspy.InputField(prefix='Topic of interest:', format=str)
    examples = dspy.InputField(prefix='Wiki page outlines of related topics for inspiration:\n', format=str)
    perspectives = dspy.OutputField()

class GenQnPrompt(dspy.Signature):
    """
    You are an experienced Wikipedia writer and want to edit a specific page. Besides your identity as a Wikipedia writer, you have a specific focus when researching the topic.
    Now, you are chatting with an expert to get information. Ask good questions to get more useful information.
    When you have no more question to ask, say "Thank you so much for your help!" to end the conversation.
    Please only ask one question at a time and don't ask what you have asked before. Your questions should be related to the topic you want to write.
    """

    topic = dspy.InputField(prefix='Topic you want to write: ', format=str)
    persona = dspy.InputField(prefix='Your specific perspective: ', format=str)
    conv = dspy.InputField(prefix='Conversation history:\n', format=str)
    question = dspy.OutputField()

class GenQueriesPrompt(dspy.Signature):
    """
    You want to answer the question using Google search. What do you type in the search box?
    Write the queries you will use in the following format:- query 1\n- query 2\n...
    """

    topic = dspy.InputField(prefix='Topic you are discussing about: ', format=str)
    question = dspy.InputField(prefix='Question you want to answer: ', format=str)
    queries = dspy.OutputField()

    \end{lstlisting}
    \captionof{lstlisting}{Prompts used in \system, corresponding to Line 4, 11, 19, 22 in Algorithm~\ref{code:pipeline}.}
    \label{lst:prompt1}
\end{figure*}

\begin{figure*}
    \begin{lstlisting}[language=Python]
class GenAnswerPrompt(dspy.Signature):
    """
    You are an expert who can use information effectively. You are chatting with a Wikipedia writer who wants to write a Wikipedia page on topic you know. You have gathered the related information and will now use the information to form a response.
    Make your response as informative as possible and make sure every sentence is supported by the gathered information.
    """

    topic = dspy.InputField(prefix='Topic you are discussing about:', format=str)
    conv = dspy.InputField(prefix='Question:\n', format=str)
    info = dspy.InputField(
        prefix='Gathered information:\n', format=str)
    answer = dspy.OutputField(prefix='Now give your response:\n')


class DirectGenOutlinePrompt(dspy.Signature):
    """
    Write an outline for a Wikipedia page.
    Here is the format of your writing:
        1. Use "#" Title" to indicate section title, "##" Title" to indicate subsection title, "###" Title" to indicate subsubsection title, and so on.
        2. Do not include other information.
    """

    topic = dspy.InputField(prefix="Topic you want to write: ", format=str)
    outline = dspy.OutputField(prefix="Write the Wikipedia page outline:\n")


class RefineOutlinePrompt(dspy.Signature):
    """
    Improve an outline for a Wikipedia page. You already have a draft outline that covers the general information. Now you want to improve it based on the information learned from an information-seeking conversation to make it more comprehensive.
    Here is the format of your writing:
        1. Use "#" Title" to indicate section title, "##" Title" to indicate subsection title, "###" Title" to indicate subsubsection title, and so on.
        2. Do not include other information.
    """

    topic = dspy.InputField(prefix="Topic you want to write: ", format=str)
    conv = dspy.InputField(prefix="Conversation history:\n", format=str)
    old_outline = dspy.OutputField(prefix="Current outline:\n", format=str)
    outline = dspy.OutputField(
        prefix='Write the Wikipedia page outline:\n')

    \end{lstlisting}
    \captionof{lstlisting}{Prompts used in \system (continue), corresponding to Line 24, 31, 32 in Algorithm~\ref{code:pipeline}.}
    \label{lst:prompt2}
\end{figure*}




\section{Automatic Evaluation Details}
\subsection{Soft Heading Recall}
\label{appendix:soft_heading_recall}
We calculate the soft heading recall between the multi-level headings in the generated outline, considered as the prediction $P$, and those in the human-written article, considered as the ground truth $G$. The calculation is based on the soft recall definition in~\citet{franti2023soft}. Given a set $A=\{Ai\}_{i=1}^K$, \textit{soft count} of an item is defined as the inverse of the sum of its similarity to other items in the set:
\begin{equation}
\resizebox{0.8\columnwidth}{!}{$%
\begin{gathered}
    \operatorname{count}\left(A_i\right)=\frac{1}{\sum_{j=1}^K \operatorname{Sim}\left(A_i, A_j\right)} \\
    \operatorname{Sim}\left(A_i, A_j\right) = \operatorname{cos}\left(\operatorname{embed}(A_i), \operatorname{embed}(A_j)\right),
\end{gathered}$%
}
\label{eq:soft_count}
\end{equation}
where $\operatorname{embed}(\cdot)$ in \refequ{eq:soft_count} is parameterized by \texttt{paraphrase-MiniLM-L6-v2} provided in the Sentence-Transformers library\footnote{\url{https://huggingface.co/sentence-transformers/paraphrase-MiniLM-L6-v2}}. The cardinality of $A$ is the sum of the counts of its individual items:
\begin{equation}
    \operatorname{card}(A)=\sum_{i=1}^K \operatorname{count}\left(A_i\right)
\end{equation}


The soft heading recall is calculated as
\begin{equation}
    \text{\textit{soft heading recall}} = \frac{\operatorname{card}(G \cap P)}{\operatorname{card}(G)},
\end{equation}
where the cardinality of intersection is defined via the union as follows:
\begin{equation}
\resizebox{0.75\columnwidth}{!}{$%
\begin{gathered}
    \operatorname{card}(G \cap P)=\\ \operatorname{card}(G)+\operatorname{card}(P)-\operatorname{card}(G \cup P).
\end{gathered}$%
}
\end{equation}




\subsection{LLM Evaluator}
\label{appendix:rubric}

\begin{table*}
\centering
\resizebox{\textwidth}{!}{%
\begin{tabular}{ll} 
\toprule
Criteria Description & \textbf{Interest Level}: How engaging and thought-provoking is the article?                                                                                                               \\
Score 1 Description  & Not engaging at all; no attempt to capture the reader's attention.                                                                                                                        \\
Score 2 Description  & Fairly engaging with a basic narrative but lacking depth.                                                                                                                                 \\
Score 3 Description  & Moderately engaging with several interesting points.                                                                                                 \\
Score 4 Description  & Quite engaging with a well-structured narrative and noteworthy points that frequently capture and retain attention.                                                                       \\
Score 5 Description  & Exceptionally engaging throughout, with a compelling narrative that consistently stimulates interest.                                                                                     \\ 
\midrule
Criteria Description & \textbf{Coherence and Organization}: Is the article well-organized and logically structured?                                                                                              \\
Score 1 Description  & Disorganized; lacks logical structure and coherence.                                                                                                                                      \\
Score 2 Description  & Fairly organized; a basic structure is present but not consistently followed.                                                                                                             \\
Score 3 Description  & Organized; a clear structure is mostly followed with some lapses in coherence.                                                                                                 \\
Score 4 Description  & Good organization; a clear structure with minor lapses in coherence.                                                                                                                      \\
Score 5 Description  & Excellently organized; the article is logically structured with seamless transitions and a clear argument.                                                                                \\ 
\midrule
Criteria Description & \textbf{Relevance and Focus}: Does the article stay on topic and maintain a clear focus?                                                                                                  \\
Score 1 Description  & Off-topic; the content does not align with the headline or core subject.                                                                                                                  \\
Score 2 Description  & Somewhat on topic but with several digressions; the core subject is evident but not consistently adhered to.                                                                              \\
Score 3 Description  & Generally on topic, despite a few unrelated details.                                                                                                                                      \\
Score 4 Description  & Mostly on topic and focused; the narrative has a consistent relevance to the core subject with infrequent digressions.                                                                    \\
Score 5 Description  & \makecell[l]{Exceptionally focused and entirely on topic; the article is tightly centered on the subject, with every piece of information contributing \\to a comprehensive understanding of the topic.}  \\ 
\midrule
Criteria Description & \textbf{Broad Coverage}: Does the article provide an in-depth exploration of the topic and have good coverage?                                                                            \\
Score 1 Description  & Severely lacking; offers little to no coverage of the topic's primary aspects, resulting in a very narrow perspective.                                                                    \\
Score 2 Description  & Partial coverage; includes some of the topic's main aspects but misses others, resulting in an incomplete portrayal.                                                                      \\
Score 3 Description  & Acceptable breadth; covers most main aspects, though it may stray into minor unnecessary details or overlook some relevant points.                                                        \\
Score 4 Description  & Good coverage; achieves broad coverage of the topic, hitting on all major points with minimal extraneous information.                                                                     \\
Score 5 Description  & \makecell[l]{Exemplary in breadth; delivers outstanding coverage, thoroughly detailing all crucial aspects of the topic without including irrelevant \\information.}                                      \\
\bottomrule
\end{tabular}
}
\caption{Scoring rubrics on a 1-5 scale for the evaluator LLM.}
\label{table:rubric}
\end{table*}


We use Prometheus\footnote{\url{https://huggingface.co/kaist-ai/prometheus-13b-v1.0}}~\citep{kim2023prometheus}, a 13B open-source evaluator LLM that can assess long-form text based on customized 1-5 scale rubric, to grade the article from the aspects of \textit{Interest level}, \textit{Coherence and Organization}, \textit{Relevance and Focus}, and \textit{Coverage}. \reftab{table:rubric} gives our grading rubric. While Prometheus is best used with a score 5 reference answer, we find adding the reference will exceed the context length limit of the model. Since \citet{kim2023prometheus} show Prometheus ratings without reference also correlate well with human preferences, we omit the reference and trim the input article to be within 2000 words by iteratively removing contents from the shortest section to ensure the input can fit into the model's context window.

\subsection{More Discussion of the Citation Quality}
\label{appendix:citation_quality_discuss}

\begin{table*}
\centering
\resizebox{\textwidth}{!}{%
\begin{tabular}{llll} 
\toprule
\textbf{Error Type}          & \textbf{Topic}                                                    & \textbf{Unsupported Sentence}                   & \textbf{Source}   \\
\midrule
Improper Inferential Linking & Lahaina, Hawaii                                                   & \makecell[l]{Throughout its history, religion has remained the\\ paramount aspect of Hawaiian life \colorbox{myyellow}{in Lahaina},\\ permeating every daily activity and significant event[5].}                                                       & \makecell[l]{[5] ``Religion, Beliefs \& Spirituality'' \\ \textit{(The source discusses religion as part of Hawaiian life}\\ \textit{but \colorbox{myyellow}{does not mention Lahania}.)}}                                      \\
\midrule[0.5pt]
Inaccurate Paraphrasing      & \makecell[l]{2022 Crimean\\Bridge explosion} & \makecell[l]{\colorbox{myyellow}{Completed in June 2020}, the bridge serves as a\\ major supply route for Russian forces in the region\\ and is significant to Russia's claim over the disputed\\ territory[2][11].} & \makecell[l]{[2] ``Crimean Bridge - Wikipedia'' \\ \textit{(The source says ``The first scheduled passenger train}\\ \textit{crossed the bridge on 25 December 2019, while the}\\ \textit{bridge was \colorbox{myyellow}{opened for freight trains on 30 June 2020}''.)}}\\
\midrule[0.5pt]
Citing Irrelevant Sources    & LK-99                                                             & \makecell[l]{For example, comparisons have been drawn between\\ the performance of LK-9 and the dynamic resolution\\ capabilities of video games such as Battlefield 2042[22].}                                                  & \makecell[l]{[22] ``Battlefield 2042 PC performance guide: The best\\ settings for a high frame rate''\\ \textit{(\colorbox{myyellow}{The source is irrelevant to LK-99.})}}  \\
\bottomrule
\end{tabular}
}
\caption{Examples of different error types of unsupported sentences.}
\label{table:example_unsupported}
\end{table*}


\begin{figure}[h]
    \centering 
    \resizebox{0.95\columnwidth}{!}{%
    \includegraphics{img/citation_error.pdf}
    }
    \caption{Error analysis of unsupported sentences in 10 sampled articles.
    }
    \label{Fig.citation_error}
\end{figure}

We use Mistral 7B-Instruct\footnote{\url{https://huggingface.co/mistralai/Mistral-7B-Instruct-v0.1}}~\citep{jiang2023mistral} to examine whether the cited passages entail the generated sentence. 
\reftab{table:citation_quality} reports the citation quality of articles produced by our approach, showing that around 15\% sentences in generated articles are unsupported by citations. We further investigate the failure cases by randomly sampling 10 articles and an author manually examines all the unsupported sentences in these articles. Besides sentences that are incorrectly split\footnote{Following~\citet{gao-etal-2023-enabling}, we check citation quality in the sentence level and split articles into sentences using NLTK \texttt{sent\_tokenize}. \texttt{sent\_tokenize} sometimes fails to split sentences correctly when the article contains special words like ``No.12847'', ``Bhatia et al.'', \etc}, lack citations, or are deemed supported by the author's judgment, our analysis identifies three main error categories (examples are given in~\reftab{table:example_unsupported}): \textit{improper inferential linking}, \textit{inaccurate paraphrasing}, and \textit{citing irrelevant sources}. 

We show the error distribution in~\reffig{Fig.citation_error}. Notably, the most common errors stem from the tendency of LLMs to form improper inferential links between different pieces of information presented in the context window. Our analysis of citation quality suggests that, in addition to avoiding hallucinations, future research in grounded text generation should also focus on preventing LLMs from making overly inferential leaps based on the provided information.



\section{Human Evaluation Details}
\label{appendix:human_eval}

We recruited 10 experienced Wikipedia editors to participate in our study by creating a research page on \texttt{Meta-Wiki}\footnote{\url{https://meta.wikimedia.org}} and reaching out to active editors who have recently approved articles for Wikipedia.\footnote{Since evaluating Wikipedia-like articles is time-consuming and requires expertise, we paid each participant 50\$ for our study.} Our participation group includes 3 editors with 1-5 years of experience, 4 with 6-10 years, and 3 with over 15 years of contribution. The study was approved by the Institutional Review Board of our institution and the participants signed the consent form through Qualtrics questionnaires before the study started.

To streamline the evaluation of grounded articles, we developed a web application, which features a side-by-side display of the article and its citation snippets, to gather ratings and open-ended feedback for each article. \reffig{Fig.ui} shows the screenshot of our web application and the full article produced by \system is included in~\reftab{table:example}. For human evaluation, we use a 1 to 7 scale for more fine-grained evaluation. The grading rubric is included in~\reftab{table:human_eval_rubric}.

We collected the pairwise preferences and the perceived usefulness of \system via an online questionnaire. Specifically, for the perceived usefulness, we request editors to rate their agreement with statements ``I think it can be specifically helpful for my pre-writing stage (\eg, collecting relevant sources, outlining, drafting).'', ``I think it will help me edit a Wikipedia article for a new topic'', ``I think it can be a potentially useful tool for the Wikipedia community'' on a Likert scale of 1-5, corresponding to \textit{Strongly disagree}, \textit{Somewhat disagree}, \textit{Neither agree nor disagree}, \textit{Somewhat agree}, \textit{Strongly agree}.

\begin{figure*}[t]
    \centering 
    \resizebox{\textwidth}{!}{%
    \includegraphics{img/ui_tool.png}
    }
    \caption{Screenshot of the web application for evaluating the grounded article.
    }
    \label{Fig.ui}
\end{figure*}

\begin{table*}
\centering
\resizebox{\textwidth}{!}{%
\begin{tabular}{l} 
\toprule
\multicolumn{1}{c}{\textbf{Interest Level}}                                                                                                                                                  \\ 
\midrule
1:~Not engaging at all; no attempt to capture the reader's attention.                                                                                                                        \\
2:~Slightly engaging with rare moments that capture attention.                                                                                                                               \\
3:~Fairly engaging with a basic narrative but lacking depth.                                                                                                                                 \\
4:~Moderately engaging with several interesting points.                                                                                                                                      \\
5:~Quite engaging with a well-structured narrative and noteworthy points that frequently capture and retain attention.                                                                       \\
6:~Very engaging with a compelling narrative that captures and mostly retains attention.                                                                                                     \\
7:~Exceptionally engaging throughout, with a compelling narrative that consistently stimulates interest.                                                                                     \\ 
\midrule
\multicolumn{1}{c}{\textbf{Coherence and Organization}}                                                                                                                                      \\ 
\midrule
1:~Disorganized; lacks logical structure and coherence.                                                                                                                                      \\
2:~Poor organization; some structure is evident but very weak.                                                                                                                               \\
3:~Fairly organized; a basic structure is present but not consistently followed.                                                                                                             \\
4:~Organized; a clear structure is mostly followed with some lapses in coherence.                                                                                                            \\
5:~Good organization; a clear structure with minor lapses in coherence.                                                                                                                      \\
6:~Very well-organized; a logical structure with transitions that effectively guide the reader.                                                                                              \\
7:~Excellently organized; the article is logically structured with seamless transitions and a clear argument.                                                                                \\ 
\midrule
\multicolumn{1}{c}{\textbf{Relevance and Focus}}                                                                                                                                             \\ 
\midrule
1:~Off-topic; the content does not align with the headline or core subject.                                                                                                                  \\
2:~Mostly off-topic with some relevant points.                                                                                                                                               \\
3:~Somewhat on topic but with several digressions; the core subject is evident but not consistently adhered to.                                                                              \\
4:~Generally on topic, despite a few unrelated details.                                                                                                                                      \\
5:~Mostly on topic and focused; the narrative has a consistent relevance to the core subject with infrequent digressions.                                                                    \\
6:~Highly relevant with a focused narrative and purpose.                                                                                                                                     \\
\makecell[l]{7:~Exceptionally focused and entirely on topic; the article is tightly centered on the subject, with every piece of information contributing to a\\ comprehensive understanding of the topic.}  \\ 
\midrule
\multicolumn{1}{c}{\textbf{Broad Coverage}}                                                                                                                                                  \\ 
\midrule
1:~Severely lacking; offers little to no coverage of the topic's primary aspects, resulting in a very narrow perspective.                                                                    \\
2:~Minimal coverage; addresses only a small selection of the topic's main aspects, with significant omissions.                                                                               \\
3:~Partial coverage; includes some of the topic's main aspects but misses others, resulting in an incomplete portrayal.                                                                      \\
4:~Acceptable breadth; covers most main aspects, though it may stray into minor unnecessary details or overlook some relevant points.                                                        \\
5:~Good coverage; achieves broad coverage of the topic, hitting on all major points with minimal extraneous information.                                                                     \\
6:~Comprehensive; provides thorough coverage of all significant aspects of the topic, with a well-balanced focus.                                                                            \\
7:~Exemplary in breadth; delivers outstanding coverage, thoroughly detailing all crucial aspects of the topic without including irrelevant information.                                      \\ 
\midrule
\multicolumn{1}{c}{\textbf{Verifiability}}                                                                                                                                                   \\ 
\midrule
1:~No supporting evidence; claims are unsubstantiated.                                                                                                                                       \\
2:~Rarely supported with evidence; many claims are unsubstantiated.                                                                                                                          \\
3:~Inconsistently verified; some claims are supported; evidence is occasionally provided.                                                                                                    \\
4:~Generally verified; claims are usually supported with evidence; however, there might be a few instances where verification is lacking                                                     \\
5:~Well-supported; claims are very well supported with credible evidence, and instances of unsupported claims are rare.                                                                      \\
6:~Very well-supported; almost every claim is substantiated with credible evidence, showing a high level of thorough verification.                                                           \\
\makecell[l]{7:~Exemplary verification; each claim is supported by robust, credible evidence from authoritative sources, reflecting strict adherence to the no \\original research policy.}                  \\
\bottomrule
\end{tabular}
}
\caption{Scoring rubrics on a 1-7 scale for human evaluation.}
\label{table:human_eval_rubric}
\end{table*}



\section{Error Analysis}
\label{appendix:error_analysis}

\begin{table*}
\centering
\resizebox{\textwidth}{!}{%
\begin{tabular}{ccl} 
\toprule
\textbf{Issue}                                    & \textbf{Mentioned Time}      & \multicolumn{1}{c}{\textbf{Example Comments}}                                                                                                                                                                                                                                                                                                                                \\ 
\midrule

\multirow{11}{*}{\makecell{Use of emotional words,\\ unneutral}}            & \multirow{11}{*}{12} & \makecell[l]{The word ``significant'' is used 17 times in this article. Vague and unsupported claims are\\ made about broader political importance and ``pivotal role[s]'', and is unencyclopedic. \\(comment on article \textit{Lahaina, Hawaii})}                                                                                                                               \\ 
\cmidrule{3-3}
                                                             &                     & \makecell[l]{{[}...] but they still have not fixed the issue of neutral point of view. It is also evident in this\\ article that the writer's standpoint is biased towards Taylor Swift. Other than that, it did \\a good job at summarizing key points and putting depth into this. \\(comment on article \textit{Speak Now (Taylor's Version)})}                              \\ 
\cmidrule{3-3}
                                                             &                     & \makecell[l]{``The film was also featured in an art and film festival hosted by The California Endowment,\\ highlighting the power of stories in reshaping narratives about communities.'' Yes, technically\\ the source says that, but it's a stretch to say in  Wikipedia voice and just sounds like \\non-neutral, promotional prose. (comment on article \textit{Gehraiyaan})}  \\ 
\midrule

\multirow{6}{*}{\makecell{Red herring fallacy, \\associating unrelated sources}} & \multirow{6}{*}{11} & \makecell[l]{Polling from America shouldn't be included and links to climate change shouldn't be\\ made unless explicitly connected by the source. (comment on article \textit{Typhoon Hinnamnor})}                                                                                                                                                                         \\ 
\cmidrule{3-3}
                                                             &                     & \makecell[l]{Sourcing seems mostly fine, though some aren't directly related (Ex. 39,40). \\(comment on article \textit{Gehraiyaan})}                                                                                                                                                                                                                                       \\ 
\cmidrule{3-3}
                                                             &                     & \makecell[l]{Here is a lengthy digression about KISS, not necessary because the article on the band \\should be linked to. (comment \textit{on article 2022 AFL Grand Final})}                                                                                                                                                                                              \\ 

\midrule
\multirow{5}{*}{Missing important information}               & \multirow{5}{*}{6}  & \makecell[l]{``One study, conducted by Sinéad Griffin, a physicist at the Lawrence Berkeley National\\ Laboratory, provided some analysis of LK-99's abilities using supercomputer simulations[20].''\\ This is not enough information  about the analysis, which would have been very useful in the \\article. (comment on article \textit{LK-99})}                               \\ 
\cmidrule{3-3}
                                                             &                     & \makecell[l]{Although the earthquake's immediate aftermath and response are adequately covered, there\\ could be more about the long-term socioeconomic impact and recovery processes. \\(comment on article \textit{2022 West Java earthquake})}                                                                                                                             \\ 
\midrule
\multirow{4}{*}{\makecell{Improper handling of\\ time-sensitive information}}  & \multirow{4}{*}{5}  & \makecell[l]{Words like ``now'' should be avoided in Wikipedia articles to prevent them from becoming\\ dated and phrases such as, ``as of December 2023'' should be used instead. \\(comment on article \textit{Cyclone Batsirai})}                                                                                                                                              \\ 
\cmidrule{3-3}
                                                             &                     & \makecell[l]{``as of December 13'' doesn't specify a year, and is old information \\(comment on article \textit{2022 West Java earthquake})}                                                                                                                                                                                                                                  \\ 
\midrule
\multirow{5}{*}{Section organization problem}                & \multirow{5}{*}{5}  & \makecell[l]{too many subsections in the ``Recovery and Rehabilitation'' section\\ (comment on article \textit{2022 West Java earthquake})}                                                                                                                                                                                                                                   \\ 
\cmidrule{3-3}
                                                             &                     & \makecell[l]{I do not like how the article is organized, with too many headers cluttering the article,\\ making it not as readable. Other than that, the AI did great work on the piece. \\(comment on article \textit{2022 Crimean Bridge explosion})}                                                                                                                       \\
\bottomrule
\end{tabular}
}
\caption{Summary of major issues found on articles produced by \system.}
\label{table:comment_summary}
\end{table*}


While articles produced by \system are preferred by both automatic metrics and human evaluation, experienced editors still identified multiple problems with the machine-generated articles. We analyze the free-form comments and summarize the major issues in \reftab{table:comment_summary}.

The primary issue raised is that the generated articles often contain emotional language and lack neutrality, primarily due to the source material. \system currently retrieves grounding sources from the Internet which is not neutral and contains considerable promotional content on its own. Addressing this bias in the pre-writing stage represents a valuable direction for future research. Another major issue is the red herring fallacy or the over-association of unrelated facts. Addressing this challenge calls for high-level sensemaking rather than mere fact-level verification.






\onecolumn
\renewcommand{\arraystretch}{1.5} %
\begin{longtable}{p{\textwidth}}
\toprule
\texttt{\small{\textbf{Taylor Hawkins}}} \\
\texttt{\small{Oliver Taylor Hawkins (1972-2022) was an American musician, best known as the drummer for the rock band Foo Fighters[1]. Born in Fort Walton, Texas, Hawkins' love for music was ignited at a young age, particularly after watching a Queen concert in 1982[2][3][5]. He kick-started his professional career as the drummer for Alanis Morissette’s band during her world tour for the hit album 'Jagged Little Pill'[8][9]. His talents were recognized by Dave Grohl, who invited him to join the Foo Fighters in 1997, marking the beginning of his impactful tenure with the band[7][8].}} \\
\texttt{\small{Hawkins was celebrated for his versatile drumming style, drawing influence from renowned drummers like Roger Taylor, Neil Peart, Phil Collins, Alex Van Halen, and Stewart Copeland[14]. His performances, marked by a unique energy and aggressive style of play, earned him recognition as one of the top rock drummers of his era[15]. Apart from his role in the Foo Fighters, Hawkins' passion for music saw him involved in numerous side projects and collaborations, cementing his place in the world of rock music[10].}} \\
\texttt{\small{Outside of his professional life, Hawkins was known for his resilience and dedication to his family. Despite personal struggles, including a near-fatal drug overdose in 2001, Hawkins remained committed to his musical career[4][9]. His legacy continues to inspire musicians and fans alike, as his contributions to rock music, coupled with his indomitable spirit, made him an unforgettable icon in the music industry[13].}} \\
\texttt{\small{Hawkins' sudden death in 2022 while on tour in Bogotá, Colombia, sent shockwaves through the music world[34]. Tributes poured in from around the globe, reflecting the respect and admiration Hawkins had garnered during his lifetime[21][31]. His life and career were honored at a star-studded tribute concert in Los Angeles, attesting to the enduring impact of his music[22].}} \\

\vspace{1pt}
\texttt{\small{\textbf{\# Early Life and Background}}} \\
\texttt{\small{Oliver Taylor Hawkins, known as Taylor Hawkins, was born and raised in Fort Walton, Texas[3]. His family moved to Laguna Beach, California when he was four years old[3]. He has two younger siblings, a brother named Jason, and a sister named Heather[3]. As a child, Hawkins was particularly influenced by his paternal grandmother, Josie Hawkins, who had grown up during the Great Depression and lived in Jackson, Mississippi[1].}} \\
\texttt{\small{During his high school days at Laguna Beach High School, from where he graduated in 1990, he became friends with Jon Davison, who later became the lead vocalist of the band Yes[2][3]. His interest in music was nurtured from an early age, particularly after watching a Queen concert in 1982 which inspired him to learn to play the drums[2][5]. He noted that music was a constant presence in his family home[5].}} \\
\texttt{\small{Despite facing certain hardships during his upbringing, including his mother's struggles with "demons", Hawkins pursued his musical ambitions[4]. He credits his older sister Heather for taking care of the family during difficult times[4].}} \\
\texttt{\small{His first major musical experience came from playing drums for Alanis Morissette's album, Jagged Little Pill, and accompanying her on the subsequent tour[3]. This marked the beginning of his professional career in the music industry.}} \\

\vspace{1pt}
\texttt{\small{\textbf{\# Career}}} \\
\texttt{\small{Taylor Hawkins began his professional music career playing in Alanis Morissette’s band during her 18-month world tour in support of the hit album 'Jagged Little Pill' from 1995 to 1997[8][9]. His performances not only in the tour but also in the music videos for “You Oughta Know”, “All I Really Want” and “You Learn” introduced him to the world of rock music and ultimately led to his meeting with Dave Grohl[8]. Throughout this time, Hawkins contributed significantly to the band's sound and performance, transforming the songs from their original drum loop format to a rock-band vibe that resonated with audiences[1][7].}} \\
\texttt{\small{In 1997, Hawkins was asked by Grohl to join the Foo Fighters, an invitation that he readily accepted[7][8]. At the time, Grohl thought it was a long shot to recruit Hawkins given that Morissette was at the height of her career, but Hawkins' desire to be a part of a rock band compelled him to make the move[7]. This marked the beginning of Hawkins' tenure as the drummer of the Foo Fighters, a role that he would play until his passing[6][9].}} \\
\texttt{\small{Apart from his work with Morissette and the Foo Fighters, Hawkins had an array of other musical experiences[10]. He drummed for Sass Jordan before joining Morissette’s touring band[10]. He was part of an ad hoc drum supergroup called SOS Allstars and filled the void for Coheed and Cambria’s 2007 album after their drummer Josh Eppard left the group[10]. In addition, Hawkins formed his own side project, the Coattail Riders, in 2005, through which he recorded his own music and took the project on the road, performing in small clubs despite the Foo Fighters' arena-status[7]. His son, Shane Hawkins, has since taken on his father's legacy, joining the Foo Fighters for a performance during the Boston Calling Music Festival in 2023[6].}}\\

\vspace{1pt}
\texttt{\small{\textbf{\# Musical Style and Influences}}} \\
\texttt{\small{Taylor Hawkins was a profound drummer, with his musical style and influences spreading across a wide array of rock genres[11]. Known for his passionate fandom of groups that came before him, Hawkins regularly expressed his admiration for bands like Rush, Genesis, and the Police, all of which featured some of the greatest drummers in rock history like Neil Peart, Phil Collins, and Stewart Copeland[11].}} \\
\texttt{\small{He was heavily influenced by his love for classic rock, as evidenced by his performances, where he covered songs from bands like Van Halen[11].}} \\
\texttt{\small{Hawkins drew influences from a variety of drumming styles, developing a signature style inspired by greats like Roger Taylor, Neil Peart, Phil Collins, Alex Van Halen, and Stewart Copeland[14]. This distinctive style and influence extended to his drum kit, which incorporated elements like rototoms and concert toms[14].}} \\
\texttt{\small{Beyond his influences, Hawkins had a unique energy that made him stand out as a drummer. His performances were recognized for their power, and he was known for his enthusiastic and aggressive style of play[15]. This earned him recognition as one of the top rock drummers of his time, with his passion for music living on through his performances[14].}} \\
\texttt{\small{Through his career, Hawkins left an indelible mark on rock music, through his distinct style, passion, and contributions to the music industry[13]. His love for music and dedication to his craft made him an unforgettable icon in the world of rock music[13].}} \\

\vspace{1pt}
\texttt{\small{\textbf{\# Personal Life}}} \\
\texttt{\small{Taylor Hawkins married Alison Hawkins, an American celebrity and entrepreneur, in 2005[18]. The couple had three children, Oliver, Annabelle, and Everleigh[19]. Hawkins' commitment to his family was evident; in fact, he even wrote a song for his middle child, Annabelle[9].}} \\
\texttt{\small{In his personal life, Hawkins had also struggled with drug use, which nearly claimed his life in a 2001 overdose[9][7][4]. However, he managed to overcome this challenge, and later expressed gratitude for the experience as a lesson that allowed him to realize the destructive path he was on[7].}} \\
\texttt{\small{Outside of his main role in the Foo Fighters, Hawkins also pursued various side projects including the Birds of Satan, NHC, and Chevy Metal. His motivation for such ventures was a constant drive to create and his love for music[7]. Hawkins was also known for his unabashed fanboy nature, often vocalizing his admiration for fellow musicians and his heroes[7].}} \\

\vspace{1pt}
\texttt{\small{\textbf{\# Legacy and Impact}}} \\
\texttt{\small{Taylor Hawkins was known for his raw and authentic drumming style, described as "courageous, damaged and unflinchingly authentic"[20]. His work with the Foo Fighters, as well as his various collaborations and side projects, made him a celebrated figure in rock ‘n’ roll[10].}} \\
\texttt{\small{Hawkins' death in 2022 was met with heartfelt tributes from colleagues and fans around the world. Notable tributes came from rock legends like Roger Taylor of Queen, who considered Hawkins as a kind, brilliant man and an inspirational mentor, likening his death to "losing a younger favourite brother"[21]. Similarly, Led Zeppelin’s Jimmy Page admired his technique, energy and spirited enthusiasm[21].}} \\
\texttt{\small{An LA tribute concert held in his honor included guest drummers like Lars Ulrich of Metallica, Travis Barker of blink-182, and Brad Wilk of Rage Against the Machine. Singers like Miley Cyrus and Alanis Morissette also performed at the concert[22].}} \\
\texttt{\small{Apart from his music, Taylor Hawkins also contributed to charities Music Support and MusiCares, both of which were chosen by the Hawkins family[23]. He had received numerous accolades throughout his career, including 27 Grammy nominations, of which he won 14[2]. In 2021, the Foo Fighters were inducted into the Rock and Roll Hall of Fame[9].}} \\

\vspace{1pt}
\texttt{\small{\textbf{\# Discography}}} \\
\texttt{\small{Taylor Hawkins also led a notable music career through his own side projects and collaborations[10]. Aside from his work with the Foo Fighters, Hawkins formed and fronted the band Taylor Hawkins \& The Coattail Riders, a project which originated from jamming sessions with his friend Drew Hester[10].}} \\
\vspace{1pt}
\texttt{\small{\textbf{\#\#\# Taylor Hawkins \& The Coattail Riders}}} \\
\texttt{\small{Taylor Hawkins \& The Coattail Riders, a band formed in 2004, have released three albums and their music spans genres including Hard Rock, Art Rock, and Alternative Rock[24][25][26]. The band grew from an initial casual jamming session, gradually evolving into a more formal arrangement that led to the production of record albums. Notably, these albums featured guest appearances by renowned musicians such as Dave Grohl, Queen's Brian May and Roger Taylor, The Cars' Elliot Easton, Perry Farrell, and Jon Davison, who is a school friend of Hawkins'[10].}} \\
\vspace{1pt}
\texttt{\small{\textbf{\#\#\# Red Light Fever}}} \\
\texttt{\small{Red Light Fever, released on April 19, 2010, was the band's first album[29][30]. Prior to its release, Hawkins revealed in an interview that the album had completed the recording and production stages, but its title and release date were yet to be determined[29]. Red Light Fever was recorded at the Foo Fighters' Studio 606 in California and featured guest musicians such as Brian May and Roger Taylor of Queen, Dave Grohl of Foo Fighters, and Elliot Easton of The Cars[29][30].}} \\
\vspace{1pt}
\texttt{\small{\textbf{\#\# Get the Money}}} \\
\texttt{\small{Get the Money, the third album from Taylor Hawkins \& The Coattail Riders, was released on November 8, 2019[29]. The album's first single, "Crossed the Line", released on October 15, 2019, featured Dave Grohl and Jon Davison, the frontman of Yes[29]. The music video for the single "I Really Blew It" also featured appearances from Grohl and Perry Farrell[29].}} \\

\vspace{1pt}
\texttt{\small{\textbf{\# Collaborations and Guest Appearances}}} \\
\texttt{\small{Throughout his career, Taylor Hawkins collaborated with various prominent artists and bands. The Coattail Riders' albums notably featured appearances from luminaries such as Brian May and Roger Taylor of Queen, Chrissie Hynde, Nancy Wilson of Heart, Sex Pistol Steve Jones and James Gang's Joe Walsh[28]. Hawkins also fronted another group, The Birds of Satan, which evolved from his heavy rock covers band, Chevy Metal[28].}} \\
\texttt{\small{Despite his diverse musical engagements, Hawkins always maintained a close allegiance with the Foo Fighters, which remained the center of his music life[7][28].}} \\

\vspace{1pt}
\texttt{\small{\textbf{\# Tragic Passing}}} \\
\texttt{\small{Taylor Hawkins, the esteemed drummer of the alt-rock band Foo Fighters, passed away suddenly on March 25, 2022, while on tour with his band in Bogotá, Colombia[34]. The official cause of death was cardiac arrest, though inquiries were raised concerning the presence of drugs in his system and their potential contribution to his death[33][34]. On the night of his passing, paramedics were called to the Four Seasons hotel in Bogotá due to reports of chest pain from an unnamed guest, later revealed to be Hawkins[34]. Unfortunately, resuscitation efforts were unsuccessful, and Hawkins was declared dead at the scene[34].}} \\
\texttt{\small{The news of Hawkins' sudden demise was announced on the morning of March 25th, 2022, which left the music world in shock[32]. The band confirmed the news with a short statement, expressing their devastation at the loss of Hawkins, whose "musical spirit and infectious laughter" would live on forever[32].}} \\
\texttt{\small{As a result of Hawkins' untimely passing, the band canceled their ongoing South American tour[33]. The festival stage at the Estéreo Picnic Festival, where the Foo Fighters were scheduled to perform that night, was transformed into a candlelight vigil in memory of Hawkins[33].}} \\
\vspace{1pt}
\texttt{\small{\textbf{\#\# Tributes and Remembrances}}} \\
\texttt{\small{In the wake of Hawkins' death, tributes from fans and colleagues alike poured in from around the world[21][31]. Among the many paying their respects were legendary rock and roll musicians like Roger Taylor, the drummer of Queen, who Hawkins credited with inspiring his own career behind the drum set[21]. In heartfelt social media posts, Taylor described Hawkins as an "inspirational mentor" and a "kind brilliant man"[21], while Led Zeppelin's Jimmy Page reminisced about sharing the stage with Hawkins and praised his "technique, energy and spirited enthusiasm"[21].}} \\
\texttt{\small{There were also numerous onstage tributes to Hawkins. Notably, Miley Cyrus expressed her grief and sent peaceful wishes to the Foo Fighters and the Hawkins family during a performance at Lollapalooza[31]. Similarly, Liam Gallagher of Oasis dedicated one of the band's biggest hits to Hawkins during a concert at the Royal Albert Hall in London[31].}} \\
\texttt{\small{Fans gathered outside the hotel where Hawkins died, lighting candles, leaving flowers, and singing the band's songs in his honor[31].}} \\
\texttt{\small{Hawkins' life and career were celebrated in a star-studded tribute concert in Los Angeles, which saw performances from over 50 musicians, including his former bands and colleagues from Def Leppard, Queen, and Foo Fighters[22].}} \\


\bottomrule
\caption{\system's generated article for ``Taylor Hawkins''. ``\#'', ``\#\#'' indicate the section title and subsection title respectively. Numbers in brackets indicate the cited references.}
\label{table:example}
\end{longtable}
